\documentclass[UTF8,12pt,a4paper,fontset=fandol]{ctexart}
% 支持从项目根目录或当前笔记目录加载共享样式。
\makeatletter
\def\input@path{{../}{../../}{../../../}{../../../../}}
\makeatother
\usepackage{researchnotes}

\title{各种平均数的概念与关系}
\author{}
\date{}

\begin{document}

\maketitle

\section{平均数的作用与基本约定}
\label{sec:means-introduction}

平均数是把一组数概括为一个代表值的方法。在不同问题中，“代表”的含义并不相同：有时关心总量的均分，有时关心比例的连续增长，有时关心速度、效率等比率，有时又希望较大的误差受到更高的惩罚。因此，平均数是一族根据问题结构选择的数值概括方式。本文先介绍四种常见平均数及其应用，再用幂平均说明它们的关系，随后讨论极限情形、连续对象上的平均、统计中的集中趋势指标与时间序列平滑，最后从变换尺度的角度统一理解并总结选择方法。

设 $n$ 为正整数，$x_1,x_2,\ldots,x_n$ 为实数。算术平均与平方平均允许数据为任意实数；本文中的几何平均、调和平均和任意实数阶的幂平均统一要求 $x_i>0$。讨论带权形式时，约定
\begin{equation}
\label{eq:means-weights}
  w_i\ge0,\qquad \sum_{i=1}^n w_i=1.
\end{equation}
权重 $w_i$ 表示第 $i$ 项在平均过程中的相对份额，可以来自人数、学分、路程等具体量。若原始权重为 $c_i\ge0$ 且 $\sum_i c_i>0$，则先令 $w_i=c_i/\sum_j c_j$。所有权重均为 $1/n$ 时，带权平均退化为等权平均；零权重项不参与平均，也不影响相应不等式的等号条件。

四种平均数的定义与应用只需代数知识。幂平均不等式的证明还使用导数与凸函数，后续连续平均的讨论使用定积分；有关证明会交代关键依据，可在学习相应分析知识后阅读。

\section{算术平均}
\label{sec:means-arithmetic}

\begin{definition}
对实数 $x_1,x_2,\ldots,x_n$，其\emph{算术平均数（arithmetic mean）}及加权算术平均数分别定义为
\begin{equation}
\label{eq:means-arithmetic}
  A=\frac{1}{n}\sum_{i=1}^n x_i,
  \qquad A_w=\sum_{i=1}^n w_i x_i,
\end{equation}
其中权重满足式~\eqref{eq:means-weights}。
\end{definition}

算术平均的核心思想是\emph{总量均分}。若 $x_i$ 表示每个人拥有的物品数量、每次测量的结果或若干门课程的成绩，将总和平均分给 $n$ 个对象，所得数值就是算术平均。带权形式则按各项所占份额分配影响，而不要求每项具有相同的重要性。

\begin{proposition}
对任意实数 $a,b$，数据 $ax_i+b$ 的加权算术平均为 $aA_w+b$，且原数据相对 $A_w$ 的加权偏差和为零，即
\[
  \sum_{i=1}^n w_i(ax_i+b)=aA_w+b,
  \qquad \sum_{i=1}^n w_i(x_i-A_w)=0.
\]
特别地，等权情形满足
\[
  \frac{(ax_1+b)+\cdots+(ax_n+b)}{n}=aA+b,
  \qquad \sum_{i=1}^n(x_i-A)=0.
\]
\end{proposition}

\begin{proof}
利用 $\sum_iw_i=1$，第一个和等于 $a\sum_iw_ix_i+b\sum_iw_i=aA_w+b$；第二个和等于 $\sum_iw_ix_i-A_w\sum_iw_i=A_w-A_w=0$。取 $w_i=1/n$，并将偏差恒等式乘以 $n$，便得到等权结论。
\end{proof}

\begin{example}
三次测量结果为 $2,4,9$；另有一名同学连续三天的通勤时间为 $28,32,30$ 分钟。分别求两组数据的算术平均，并解释后者的含义。
\end{example}

\begin{solve}
两组数据的算术平均分别为
\[
  \frac{2+4+9}{3}=5,
  \qquad \frac{28+32+30}{3}=30\text{ 分钟}.
\]
前一个结果说明平均数未必等于任何一次观测值；后一个结果表示将三天总共 $90$ 分钟的通勤时间均分到每天，每天相当于 $30$ 分钟。
\end{solve}

\begin{example}
某门课程的平时作业、期中考试、期末考试权重依次为 $20\%,30\%,50\%$，对应成绩为 $80,90,86$，求总评成绩。若另有三门课的学分依次为 $2,3,5$，成绩同样为 $80,90,86$，求按学分加权的平均成绩。
\end{example}

\begin{solve}
第一种情形的权重已经归一化，总评为
\[
  0.2\cdot80+0.3\cdot90+0.5\cdot86=86.
\]
第二种情形的总学分为 $2+3+5=10$，归一化权重仍为 $0.2,0.3,0.5$，因此平均成绩也为 $86$。两者的数学形式相同，但权重分别表示考核比例与学分比例。在其他成绩不变时，第一种情形的期末成绩每提高 $1$ 分，总评提高 $0.5$ 分，这说明权重如何决定一项数据的影响。
\end{solve}

\section{几何平均}
\label{sec:means-geometric}

\begin{definition}
对正数 $x_1,x_2,\ldots,x_n$，其\emph{几何平均数（geometric mean）}及加权几何平均数分别定义为
\begin{equation}
\label{eq:means-geometric}
  G=\left(\prod_{i=1}^n x_i\right)^{1/n},
  \qquad G_w=\prod_{i=1}^n x_i^{w_i}.
\end{equation}
\end{definition}

几何平均的核心思想是\emph{乘法结构的均分}。一个数量依次变为前一步的 $x_1,x_2,\ldots,x_n$ 倍，最终就变为初始值的 $x_1x_2\cdots x_n$ 倍。若将每一步都改为同一个倍率 $G$，同时要求 $n$ 步后的总效果保持不变，就必须满足
\[
  G^n=x_1x_2\cdots x_n.
\]
因此，几何平均给出等效的平均倍率；对于增长率 $r_i$，应先转换为正的增长倍率 $1+r_i$，再求几何平均，最后减去 $1$。

由对数的乘法法则，
\begin{equation}
\label{eq:means-geometric-log}
  \ln G=\frac{1}{n}\sum_{i=1}^n\ln x_i,
  \qquad \ln G_w=\sum_{i=1}^n w_i\ln x_i.
\end{equation}
这表明几何平均也可以理解为：先用对数将相乘的倍率转化为相加的量，再求算术平均，最后指数还原。这个视角解释了几何平均为何适合连续按比例变化的过程；若各阶段时长不同，计算单位时间的等效增长率时，还需要将倍率与时长对应起来。

\begin{example}
某投资第一年增长 $20\%$，第二年下降 $10\%$，第三年增长 $50\%$。求保持三年累计效果不变的等效年均增长率。
\end{example}

\begin{solve}
三年的增长倍率为 $1.2,0.9,1.5$，累计倍率为 $1.2\cdot0.9\cdot1.5=1.62$。因此等效年均增长倍率与增长率分别为
\[
  G=\sqrt[3]{1.62},
  \qquad r=\sqrt[3]{1.62}-1.
\]
这里每一年均以前一年的结果为基数，故不能直接取 $20\%,-10\%,50\%$ 的算术平均作为等效年均增长率。
\end{solve}

\section{调和平均}
\label{sec:means-harmonic}

\begin{definition}
对正数 $x_1,x_2,\ldots,x_n$，其\emph{调和平均数（harmonic mean）}及加权调和平均数分别定义为
\begin{equation}
\label{eq:means-harmonic}
  H=\frac{n}{\sum_{i=1}^n1/x_i},
  \qquad H_w=\frac{1}{\sum_{i=1}^n w_i/x_i}.
\end{equation}
\end{definition}

调和平均的核心思想是\emph{倒数结构的均分}。设第 $i$ 段路程为 $d_i>0$、速度为 $x_i>0$，则耗时为 $d_i/x_i$，全程平均速度为
\[
  \frac{\sum_i d_i}{\sum_i d_i/x_i}
  =\frac{1}{\sum_i w_i/x_i},
  \qquad w_i=\frac{d_i}{\sum_jd_j}.
\]
因此，按路程份额加权时得到调和平均，等路程时得到等权调和平均。相反，若已知各段时间 $t_i>0$，则平均速度为 $\sum_it_ix_i/\sum_it_i$，是按时间份额加权的算术平均。

这一推导同样适用于“单位资源产生多少结果”等比率：若 $x_i=a_i/b_i$，其中 $a_i,b_i>0$，则总比率 $\sum_i a_i/\sum_i b_i$ 按分子份额加权时是调和平均，按分母份额加权时是算术平均。\emph[caution]{仅知道数据是速度或效率，还不足以决定使用调和平均；必须先确定哪些量可以相加，以及权重代表什么。}

\begin{example}
一辆车去程速度为 $60$ 千米/小时，返程速度为 $40$ 千米/小时，往返路程相同且途中无停留。求全程平均速度。
\end{example}

\begin{solve}
设单程路程为 $d>0$ 千米，则总路程为 $2d$ 千米，总时间为 $d/60+d/40$ 小时。因此
\[
  \overline v=\frac{2d}{d/60+d/40}
  =\frac{2}{1/60+1/40}=48\text{ 千米/小时}.
\]
算术平均 $(60+40)/2=50$ 千米/小时对应的是两种速度持续时间相同的情形；本题固定的是相同路程，低速返程耗时更长，故结果为 $48$ 千米/小时。
\end{solve}

\section{平方平均与误差大小}
\label{sec:means-quadratic}

\begin{definition}
对实数 $x_1,x_2,\ldots,x_n$，其\emph{平方平均数（quadratic mean）}，又称\emph{均方根（root mean square, RMS）}，及其带权形式分别定义为
\begin{equation}
\label{eq:means-quadratic}
  Q=\sqrt{\frac{1}{n}\sum_{i=1}^n x_i^2},
  \qquad Q_w=\sqrt{\sum_{i=1}^n w_i x_i^2}.
\end{equation}
\end{definition}

平方平均先平方、再求算术平均、最后开平方：平方使正负号不再互相抵消，开方则将结果还原到原来的量纲。它衡量数据相对零点的整体大小；用于有符号误差时，可以反映误差强度。平方会放大绝对值较大的观测值的影响，因此平方平均对大误差更敏感。它还与交流电有效值、信号强度和向量长度有关；例如向量 $(x_1,\ldots,x_n)$ 的欧几里得长度为 $\sqrt n\,Q$。

\begin{example}
某工厂零件的设计长度为 $100$ 毫米，三个零件相对设计长度的误差为 $-2,0,2$ 毫米。另一次测量的三个误差为 $-1,2,-2$ 毫米。分别计算两组误差的算术平均与均方根，说明它们反映的信息。
\end{example}

\begin{solve}
第一组的平均误差为 $(-2+0+2)/3=0$，均方根为
\[
  Q=\sqrt{\frac{(-2)^2+0^2+2^2}{3}}
  =\sqrt{\frac83}\text{ 毫米}\approx1.63\text{ 毫米}.
\]
平均误差为零说明这三个零件的正负偏差恰好抵消，并不表示每个零件都没有误差；$1.63$ 毫米描述的是均方根意义下的误差大小。第二组满足
\[
  A=\frac{-1+2-2}{3}=-\frac13\text{ 毫米},
  \qquad Q=\sqrt{\frac{(-1)^2+2^2+(-2)^2}{3}}
  =\sqrt3\text{ 毫米}.
\]
这里 $A$ 反映有方向的平均偏差，$Q$ 反映相对零误差的整体偏离程度，两者回答不同的问题。
\end{solve}

\begin{proposition}
\label{prop:means-rms-decomposition}
令 $V_w=\sum_iw_i(x_i-A_w)^2$ 为数据相对加权均值的加权平均平方偏差，则
\begin{equation}
\label{eq:means-rms-decomposition}
  Q_w^2=A_w^2+V_w.
\end{equation}
因此 $Q_w\ge|A_w|$，等号成立当且仅当所有正权重对应的数据相等。
\end{proposition}

\begin{proof}
展开平方并使用 $\sum_iw_ix_i=A_w$ 与 $\sum_iw_i=1$，得到
\[
  V_w=\sum_iw_ix_i^2-2A_w\sum_iw_ix_i+A_w^2\sum_iw_i
  =Q_w^2-A_w^2.
\]
又因为每一项 $w_i(x_i-A_w)^2$ 非负，$V_w=0$ 当且仅当所有满足 $w_i>0$ 的 $x_i$ 均等于 $A_w$，故结论成立。
\end{proof}

式~\eqref{eq:means-rms-decomposition} 区分了相对零点的大小与相对均值的离散程度。等权情形下，$\sqrt{n^{-1}\sum_i(x_i-A)^2}$ 是以 $n$ 为分母的数据标准差，只有 $A=0$ 时才等于原数据的均方根；统计推断中常用的样本标准差在 $n\ge2$ 时以 $n-1$ 为分母，不能与这里的量混同。

\section{幂平均及其不等式}
\label{sec:means-power}

\begin{definition}
对正数 $x_1,x_2,\ldots,x_n$ 和实数 $p$，\emph{$p$ 阶幂平均（power mean）}定义为
\begin{equation}
\label{eq:means-power}
  M_p=
  \begin{cases}
    \left(\dfrac{1}{n}\sum_{i=1}^n x_i^p\right)^{1/p}, & p\ne0,\\[1.2em]
    \left(\prod_{i=1}^n x_i\right)^{1/n}, & p=0.
  \end{cases}
\end{equation}
相应的加权幂平均为
\begin{equation}
\label{eq:means-weighted-power}
  M_p(w)=
  \begin{cases}
    \left(\displaystyle\sum_{i=1}^n w_i x_i^p\right)^{1/p}, & p\ne0,\\[1.2em]
    \displaystyle\prod_{i=1}^n x_i^{w_i}, & p=0.
  \end{cases}
\end{equation}
\end{definition}

幂平均把前述四种平均数统一到同一个框架中：
\begin{equation}
\label{eq:means-power-special-cases}
  M_{-1}=H,\qquad M_0=G,\qquad M_1=A,\qquad M_2=Q.
\end{equation}
带权情形也有 $M_{-1}(w)=H_w$、$M_0(w)=G_w$、$M_1(w)=A_w$、$M_2(w)=Q_w$。这些统一关系在正数范围内使用；算术平均和平方平均本身允许负数，但任意实数阶幂平均的上述定义不作这种扩展。随着 $p$ 增大，平均值向较大的数据一侧移动；随着 $p$ 减小，平均值向较小的数据一侧移动。下面给出这一说法的严格结论，并说明 $p=0$ 的定义与极限一致。

\begin{theorem}[幂平均不等式]
\label{thm:means-power-inequality}
设 $x_i>0$，权重满足式~\eqref{eq:means-weights}。则 $p\mapsto M_p(w)$ 在实数轴上连续；对任意实数 $p<q$，有
\begin{equation}
\label{eq:means-power-inequality}
  M_p(w)\le M_q(w).
\end{equation}
等号成立当且仅当所有满足 $w_i>0$ 的 $x_i$ 相等。特别地，等权情形有 $M_p\le M_q$，且等号成立当且仅当 $x_1=\cdots=x_n$。
\end{theorem}

\begin{proof}
\emph[tip]{对幂平均取对数，将比较不同阶数的问题转化为比较同一凸函数的割线斜率。}令 $y_i=\ln x_i$，并定义
\[
  \varphi(t)=\ln\left(\sum_{i=1}^n w_i e^{t y_i}\right),
  \qquad
  \alpha_i(t)=\frac{w_i e^{t y_i}}{\sum_{j=1}^n w_j e^{t y_j}}.
\]
因为至少有一个权重为正，分母对所有实数 $t$ 都为正，所以 $\varphi$ 在实数轴上二阶连续可导，且 $\alpha_i(t)\ge0$、$\sum_i\alpha_i(t)=1$。求导得到
\[
  \varphi'(t)=\sum_i\alpha_i(t)y_i,
  \qquad
  \varphi''(t)=\sum_i\alpha_i(t)y_i^2
  -\left(\sum_i\alpha_i(t)y_i\right)^2.
\]
记 $\mu(t)=\sum_i\alpha_i(t)y_i$，则
\[
  \varphi''(t)=\sum_i\alpha_i(t)\bigl(y_i-\mu(t)\bigr)^2\ge0.
\]
故 $\varphi$ 为凸函数。这里凸性是指对任意实数 $u,v$ 与 $\theta\in[0,1]$，有
\[
  \varphi((1-\theta)u+\theta v)
  \le(1-\theta)\varphi(u)+\theta\varphi(v).
\]
二阶导数非负保证导数不减，由中值定理比较相邻区间的割线斜率即可得到这一不等式。

又由权重和为 $1$，有 $\varphi(0)=0$。对 $t\ne0$，
\[
  \ln M_t(w)=\frac{\varphi(t)}{t}.
\]
由导数定义，
\[
  \lim_{t\to0}\frac{\varphi(t)}{t}
  =\varphi'(0)=\sum_iw_i\ln x_i=\ln M_0(w).
\]
因此 $\ln M_t(w)$ 在 $t=0$ 处连续，其余各点的连续性由公式直接得到；再利用指数函数的连续性，即知 $M_t(w)$ 在整个实数轴上连续。

为比较包括负阶数在内的情形，先注意：对任意 $u<v<z$，将 $v$ 写成 $u,z$ 的凸组合，并整理凸性不等式，可得
\begin{equation}
\label{eq:means-convex-secants}
  \frac{\varphi(v)-\varphi(u)}{v-u}
  \le\frac{\varphi(z)-\varphi(u)}{z-u}
  \le\frac{\varphi(z)-\varphi(v)}{z-v}.
\end{equation}
当 $0<p<q$ 时，取 $(u,v,z)=(0,p,q)$ 并使用第一个不等式；当 $p<q<0$ 时，取 $(u,v,z)=(p,q,0)$ 并使用第二个不等式；当 $p<0<q$ 时，取 $(u,v,z)=(p,0,q)$ 并比较最左端与最右端。三种情形均给出
\[
  \frac{\varphi(p)}{p}\le\frac{\varphi(q)}{q}.
\]
若 $p=0$ 或 $q=0$，则由刚证明的连续性取极限得到同样结论。于是 $\ln M_p(w)\le\ln M_q(w)$，再由指数函数严格递增得到式~\eqref{eq:means-power-inequality}。

最后讨论等号。若正权重对应的数据不全相等，则相应的 $y_i$ 不全相等，且对应的 $\alpha_i(t)$ 始终为正，因此 $\varphi''(t)>0$ 对所有 $t$ 成立。这时 $\varphi$ 严格凸，式~\eqref{eq:means-convex-secants} 中的两个不等式均严格，故上述非零阶数之间的比较均严格。若一个阶数为零，由中值定理，$\varphi(t)/t=\varphi'(\xi)$，其中 $\xi$ 严格位于 $0$ 与 $t$ 之间；由于 $\varphi'$ 严格递增，该值在 $t<0$ 时小于 $\varphi'(0)$，在 $t>0$ 时大于 $\varphi'(0)$，仍得到严格不等式。反之，若所有正权重对应的数据均为同一正数 $c$，则每一阶幂平均都等于 $c$，故等号条件成立。
\end{proof}

\begin{corollary}
\label{cor:means-chain}
对正数 $x_1,\ldots,x_n$，有
\begin{equation}
\label{eq:means-chain}
  H\le G\le A\le Q.
\end{equation}
任意相邻两项取等号，当且仅当 $x_1=\cdots=x_n$。带权情形相应满足 $H_w\le G_w\le A_w\le Q_w$，等号条件为所有正权重对应的数据相等。
\end{corollary}

\begin{proof}
在定理~\ref{thm:means-power-inequality} 中依次比较阶数 $-1,0,1,2$，再使用式~\eqref{eq:means-power-special-cases} 及其带权形式即可。
\end{proof}

这条不等式链比较的是同一组正数据、同一组权重下的平均值，数值较大并不表示某种平均数在实际问题中更合适。对于允许负数的数据，命题~\ref{prop:means-rms-decomposition} 中的 $Q_w\ge|A_w|$ 仍成立，但不能直接套用整条正数平均值不等式链。

\section{最值平均与极限情形}
\label{sec:means-extreme-limits}

\begin{proposition}
\label{prop:means-extreme-limits}
对正数 $x_1,\ldots,x_n$，有
\[
  \lim_{p\to+\infty}M_p=\max_{1\le i\le n}x_i,
  \qquad
  \lim_{p\to-\infty}M_p=\min_{1\le i\le n}x_i.
\]
带权情形的两个极限分别为正权重对应的数据中的最大值与最小值。
\end{proposition}

\begin{proof}
令 $I_+=\{i:w_i>0\}$，并取 $j\in I_+$ 使 $x_j=b=\max_{i\in I_+}x_i$。对 $p>0$，有
\[
  w_jb^p\le\sum_iw_ix_i^p\le b^p,
  \qquad w_j^{1/p}b\le M_p(w)\le b.
\]
由于 $w_j>0$ 且 $w_j^{1/p}\to1$，夹逼定理给出 $\lim_{p\to+\infty}M_p(w)=b$。若用 $M_r(w;x)$ 表示数据 $x_1,\ldots,x_n$ 的加权幂平均，则对 $r>0$，
\[
  M_{-r}(w;x)=\frac{1}{M_r(w;1/x)},
\]
其中 $1/x$ 表示倒数数据 $1/x_1,\ldots,1/x_n$。对倒数数据使用正无穷处的结论，再取倒数，得到负无穷处的极限为 $\min_{i\in I_+}x_i$。取 $w_i=1/n$ 即得等权结论。
\end{proof}

因此，最大值与最小值可以视为幂平均的两个极限端点。结合定理~\ref{thm:means-power-inequality}，幂平均构成从最小值向最大值连续、单调变化的一族平均数。直观上，大的正幂使最大项主导求和；大的负幂使最小正数对应的项主导求和，再通过负指数还原。

\section{积分平均}
\label{sec:means-integral}

\begin{definition}
设 $a<b$，实函数 $f$ 在 $[a,b]$ 上黎曼可积，则其\emph{积分平均（integral mean）}定义为
\begin{equation}
\label{eq:means-integral}
  \overline f=\frac{1}{b-a}\int_a^b f(x)\,\mathrm dx.
\end{equation}
\end{definition}

积分平均是连续对象上的算术平均：离散情形将有限多个数相加再除以个数，连续情形则以积分替代求和，以区间长度替代个数。例如，物体沿固定方向运动且不折返时，在时间区间 $[a,b]$ 上对速度积分得到路程，再除以总时间便得到这一时间段的平均速度。

\begin{proposition}[积分平均值定理]
若 $a<b$，$f$ 在 $[a,b]$ 上连续，则存在 $\xi\in[a,b]$，使得
\[
  f(\xi)=\frac{1}{b-a}\int_a^b f(x)\,\mathrm dx.
\]
\end{proposition}

\begin{proof}
由连续函数在闭区间上取得最值，记其最小值、最大值为 $m,M$。对 $m\le f(x)\le M$ 积分并除以 $b-a>0$，得到 $m\le\overline f\le M$。再由介值定理，$f$ 在区间内某点取到 $\overline f$。
\end{proof}

这一结论说明连续函数的平均值至少在某一点被实际取到。若仅假设可积，则仍可定义平均值，却不保证函数取到该值；例如在 $[0,1/2)$ 取 $0$、在 $[1/2,1]$ 取 $1$ 的函数，其积分平均为 $1/2$，但函数值只有 $0$ 和 $1$。

\section{对数平均}
\label{sec:means-logarithmic}

\begin{definition}
对正数 $a,b$，其\emph{对数平均（logarithmic mean）}定义为
\begin{equation}
\label{eq:means-logarithmic}
  L(a,b)=
  \begin{cases}
    \dfrac{b-a}{\ln b-\ln a},&a\ne b,\\
    a,&a=b.
  \end{cases}
\end{equation}
\end{definition}

当 $0<a<b$ 时，利用式~\eqref{eq:means-integral}，可将对数平均理解为函数 $x\mapsto1/x$ 在 $[a,b]$ 上积分平均的倒数，因为
\[
  \frac{1}{b-a}\int_a^b\frac{1}{x}\,\mathrm dx
  =\frac{\ln b-\ln a}{b-a},
  \qquad
  L(a,b)=\left(\frac{1}{b-a}\int_a^b\frac{1}{x}\,\mathrm dx\right)^{-1}.
\]
若 $a>b$，交换两端点即可，所得 $L(a,b)$ 不变。另一种表示是
\begin{equation}
\label{eq:means-logarithmic-integral}
  L(a,b)=\int_0^1 a^{1-t}b^t\,\mathrm dt.
\end{equation}
当 $a\ne b$ 时，对 $a\exp(t\ln(b/a))$ 直接积分即可得到该式；当 $a=b$ 时，被积函数恒为 $a$。因此，对数平均也表示一个量从 $a$ 到 $b$ 按指数方式变化时，在归一化时间区间上的算术平均，这说明它与指数增长及相应的传热平均温差计算的联系。由原定义还可得
\[
  \lim_{b\to a}L(a,b)
  =\left(\lim_{b\to a}\frac{\ln b-\ln a}{b-a}\right)^{-1}=a,
\]
故 $a=b$ 处的定义与极限一致。

\begin{proposition}
对正数 $a,b$，记 $G(a,b)=\sqrt{ab}$、$A(a,b)=(a+b)/2$，则
\[
  G(a,b)\le L(a,b)\le A(a,b),
\]
且任意一个不等式取等号，当且仅当 $a=b$。
\end{proposition}

\begin{proof}
由推论~\ref{cor:means-chain} 的加权几何平均不超过加权算术平均，对 $t\in[0,1]$ 有
\[
  a^{1-t}b^t\le(1-t)a+tb.
\]
在 $[0,1]$ 上积分，并使用式~\eqref{eq:means-logarithmic-integral}，得到 $L(a,b)\le(a+b)/2$。另一方面，两个正数 $a^{1-t}b^t$ 与 $a^tb^{1-t}$ 的乘积为 $ab$，故等权算术--几何平均不等式给出
\[
  a^{1-t}b^t+a^tb^{1-t}\ge2\sqrt{ab}.
\]
两边积分，左侧的两个积分经变量代换 $t\mapsto1-t$ 可知相等，于是 $2L(a,b)\ge2\sqrt{ab}$。若 $a\ne b$，上界中的不等式对每个 $0<t<1$ 均严格，下界中的不等式除 $t=1/2$ 外均严格；被积函数连续，故积分后的两个不等式均严格。若 $a=b$，三个平均数都等于 $a$。
\end{proof}

\section{统计中的集中趋势指标}
\label{sec:means-statistical-location}

在统计中，算术平均常用于描述数据的集中趋势；若数据含有异常值或明显偏斜，也会使用中位数、众数和截尾平均等\emph{集中趋势指标（measures of central tendency）}。它们所概括的信息不同，并非都属于幂平均或由同一种代数运算产生。

\subsection{中位数}

\begin{definition}
将实数数据从小到大排列为 $x_{(1)}\le\cdots\le x_{(n)}$，其\emph{中位数（median）}定义为
\[
  \operatorname{Med}(x)=
  \begin{cases}
    x_{((n+1)/2)},&n\text{ 为奇数},\\
    \dfrac{x_{(n/2)}+x_{(n/2+1)}}{2},&n\text{ 为偶数}.
  \end{cases}
\]
\end{definition}

中位数由排序后居中的观测值决定。只要这些居中位置的值保持不变，将尾部极端值进一步增大或减小就不会改变中位数，因此它比算术平均更能抵抗少量极端值的影响。偶数样本量时，本文采用中间两个数的算术平均这一常用约定。

\begin{example}
比较数据 $2,3,4,5,100$ 的算术平均与中位数，并说明差异的来源。
\end{example}

\begin{solve}
这组数据的算术平均为 $(2+3+4+5+100)/5=22.8$，中位数为 $4$。较大的观测值 $100$ 显著拉高了算术平均；若关心排序居中的典型水平，中位数更贴近该目标，若关心总量均分，则仍应使用算术平均。
\end{solve}

\subsection{众数}

\begin{definition}
有限数据中出现次数最多的取值称为\emph{众数（mode）}；若多个取值的出现次数并列最多，则它们均为众数。
\end{definition}

众数适合描述类别型或离散数据的最常见取值，例如服装尺码、调查选项或离散考试分数。按上述约定，若每个取值出现次数都相同，则它们都达到最高频数，此时众数无法突出某个典型取值。对于已经分组的数值数据，可报告出现次数最多的组，但不能据此确定原始数据的精确众数。

\subsection{截尾平均}

\begin{definition}
设 $k$ 为整数且 $0\le2k<n$。将数据排序后，删除最小的 $k$ 项与最大的 $k$ 项，剩余数据的算术平均称为相应的\emph{截尾平均（trimmed mean）}，即
\[
  T_k=\frac{1}{n-2k}\sum_{i=k+1}^{n-k}x_{(i)}.
\]
\end{definition}

若按每端比例 $\beta\in[0,1/2)$ 截尾，本文取 $k=\lfloor\beta n\rfloor$，其中 $\lfloor\cdot\rfloor$ 表示向下取整。例如每端 $10\%$ 截尾就是分别删除两端各 $\lfloor0.1n\rfloor$ 项，再对中间部分求平均；当 $n$ 是 $10$ 的倍数时，每端恰好删除 $10\%$ 的观测值。截尾平均保留了中间大部分数值信息，同时减弱尾部极端值的影响，在数据概括方式上介于使用全部数值的算术平均与主要依赖居中位置的中位数之间。

\section{移动平均}
\label{sec:means-moving}

\begin{definition}
对实数时间序列 $x_1,x_2,\ldots,x_N$，设 $N$ 为正整数，窗口长度 $k$ 为满足 $1\le k\le N$ 的整数。长度为 $k$ 的\emph{简单移动平均（simple moving average）}定义为
\[
  m_t=\frac{1}{k}\sum_{j=t-k+1}^{t}x_j,
  \qquad t=k,k+1,\ldots,N.
\]
\end{definition}

移动平均用滑动窗口反复计算局部平均，可用于减弱短期波动、观察较缓慢的变化。它输出一串随时间变化的数，而非整组数据的一个代表值；窗口越长，使用的历史范围越大，但对新变化的响应通常也越滞后。

\begin{definition}
给定参数 $0<\alpha<1$，以 $s_1=x_1$ 初始化，递推定义
\begin{equation}
\label{eq:means-exponential-moving}
  s_t=\alpha x_t+(1-\alpha)s_{t-1},
  \qquad t=2,\ldots,N,
\end{equation}
所得序列称为\emph{指数移动平均（exponential moving average）}。
\end{definition}

反复展开式~\eqref{eq:means-exponential-moving}，得到
\[
  s_t=(1-\alpha)^{t-1}x_1
  +\alpha\sum_{j=2}^{t}(1-\alpha)^{t-j}x_j,
  \qquad t\ge2.
\]
各权重非负且和为 $1$，其中后续观测 $x_j$ 的权重随其距当前时刻的间隔指数衰减，$x_1$ 的系数则保留了初始化的影响。参数 $\alpha$ 越大，当前观测在更新中占比越高；$\alpha$ 越小，历史状态占比越高，平滑通常更强，但响应也更慢。

\section{拟算术平均：从变换尺度统一理解}
\label{sec:means-quasi-arithmetic}

许多平均数可以通过同一过程构造：先变换数据，在变换后的尺度上求算术平均，再用反函数还原。为了保证中间的平均值落在反函数的定义域内，需要明确变换函数的条件。

\begin{definition}
设 $I\subseteq\mathbb R$ 为区间，$\varphi:I\to\mathbb R$ 连续且严格单调，$x_1,\ldots,x_n\in I$。由 $\varphi$ 生成的\emph{拟算术平均（quasi-arithmetic mean）}及其带权形式定义为
\begin{equation}
\label{eq:means-quasi-arithmetic}
  M_{\varphi}=\varphi^{-1}\left(\frac{1}{n}\sum_{i=1}^n\varphi(x_i)\right),
  \qquad
  M_{\varphi}(w)=\varphi^{-1}\left(\sum_{i=1}^n w_i\varphi(x_i)\right).
\end{equation}
\end{definition}

连续性保证 $\varphi(I)$ 仍为区间，故其中有限个值的加权平均仍属于 $\varphi(I)$；严格单调性保证反函数在该像集上存在，所以式~\eqref{eq:means-quasi-arithmetic} 有定义。只假设严格单调而不假设连续，像集可能存在间隙，便不能保证这个构造对任意数据都成立。

取不同的生成函数，可得到
\[
\begin{aligned}
  \varphi(x)=x,\quad I=\mathbb R
  &\quad\Longrightarrow\quad M_{\varphi}=A,\\
  \varphi(x)=\ln x,\quad I=(0,+\infty)
  &\quad\Longrightarrow\quad M_{\varphi}=G,\\
  \varphi(x)=1/x,\quad I=(0,+\infty)
  &\quad\Longrightarrow\quad M_{\varphi}=H,\\
  \varphi(x)=x^2,\quad I=[0,+\infty)
  &\quad\Longrightarrow\quad M_{\varphi}=Q.
\end{aligned}
\]
特别地，$x^2$ 在整个实数轴上不严格单调，所以这一拟算术平均表示限于非负数据；平方平均本身仍可按式~\eqref{eq:means-quadratic} 对任意实数数据计算。对正数数据，取 $\varphi(x)=x^p$（$p\ne0$）可得到 $M_p$，取 $\varphi(x)=\ln x$ 可得到 $M_0$，带权形式也相应成立。

这一观点把平均数的选择与问题中的变换尺度联系起来：在普通数量尺度上均分得到算术平均，在对数尺度上均分得到几何平均，在倒数尺度上均分得到调和平均，在平方尺度上均分并取非负平方根得到平方平均。它统一了这些平均数的构造，但并不要求中位数、众数等其他概括量也具有同样的表示。

\section{如何选择平均数}
\label{sec:means-selection}

选择平均数，首先要明确希望保留的总效果与权重的来源。\emph{总量可加}且关心均分时，用算术平均；若各项代表不同人数、学分或持续时间，则按相应份额加权。\emph{倍率相乘}且关心累计效果时，用几何平均，例如复利和连续增长问题，应先把增长率转换为倍率，再计算等效增长率。

若可加总量与观测比率的倒数成正比，例如按给定路程计算各段耗时，则使用具有相应权重的调和平均。\emph[caution]{比率的分子、分母以及加权依据必须一致。}第~\ref{sec:means-harmonic} 节中的等路程与等时间速度问题，正好说明同一组速度可以因条件不同而需要不同的平均方式。

若关心有符号数据相对零点的大小，并希望较大的绝对误差产生更大影响，可用平方平均；若关心围绕均值的波动，则应先中心化再求均方根，二者的联系见式~\eqref{eq:means-rms-decomposition}。若关心含少量极端值的数据的典型位置，可考虑中位数或截尾平均；若关心最常见的类别，可用众数。时间序列的局部平滑则应选择移动平均，并明确窗口长度或递推参数。一个概括量是否合适，取决于它回答的问题，而非仅取决于它在不等式链中的大小。

\section{常见平均数对照表}
\label{sec:means-comparison}

表~\ref{tab:means-comparison} 汇总主要公式及其使用条件。表中的几何平均、调和平均和幂平均均采用 $x_i>0$ 的约定；算术平均和平方平均允许实数数据，各带权形式见前文定义。

\begin{center}
\captionof{table}{常见平均数与集中趋势指标}
\label{tab:means-comparison}
\begin{tabularx}{\textwidth}{L{3.1cm}Y Y}
\toprule
名称 & 公式或定义 & 适用直观与条件 \\
\midrule
算术平均 & $\displaystyle A=\frac{1}{n}\sum_{i=1}^n x_i$ & 总量均分，适合可加量 \\
几何平均 & $\displaystyle G=\left(\prod_{i=1}^n x_i\right)^{1/n}$ & 保持累计倍率，增长率先转为倍率 \\
调和平均 & $\displaystyle H=\frac{n}{\sum_{i=1}^n1/x_i}$ & 平均倒数再取倒数，例如等路程速度 \\
平方平均 & $\displaystyle Q=\sqrt{\frac{1}{n}\sum_{i=1}^n x_i^2}$ & 相对零点的大小，较重惩罚大误差 \\
幂平均 & $\displaystyle M_p=\left(\frac{1}{n}\sum_{i=1}^n x_i^p\right)^{1/p}$ & 此式用于 $p\ne0$；$M_0=G$，见式~\eqref{eq:means-power} \\
对数平均 & $\displaystyle L(a,b)=\frac{b-a}{\ln b-\ln a}$ & $a,b>0$ 且 $a\ne b$；$L(a,a)=a$，联系指数变化 \\
积分平均 & $\displaystyle\overline f=\frac{1}{b-a}\int_a^b f(x)\,\mathrm dx$ & $a<b$ 且 $f$ 可积，连续对象上的均分 \\
中位数 & 排序后的中间项；偶数项时取中间两项均值 & 按排序位置概括集中趋势 \\
众数 & 出现次数最多的取值 & 描述最常见的类别或离散值，可有多个 \\
截尾平均 & 删除两端各 $k$ 项后求均值 & $0\le2k<n$，保留中间数据并减弱尾部影响 \\
移动平均 & 窗口内求均值或作指数递推 & 输出局部平均序列，需指定窗口或参数与初值 \\
\bottomrule
\end{tabularx}
\end{center}

\end{document}
